% GENERATED FILE. Edit canonical lessons, then run npm run generate:books.
% canonical-source-hash: fnv1a64:1d7c5aacf7425635
% canonical-lessons: FR-C399-panne, FR-C399-planter, FR-C399-bloquer, FR-C399-redemarrer, FR-C399-virus, FR-R399-revision-398-399, FR-R399-revision-final-399

\chapter{Ça ne marche pas}
\label{ch:fr-ca-ne-marche-pas}
\hlchaptermodality{399}

\begin{chapteropening}
\textbf{By the end of this chapter:} \emph{I can say what is wrong with a device: it has broken down, crashed or frozen, needs restarting or has a virus.}

Complete the last lesson of chapter 399: I can say what is wrong with a device: it has broken down, crashed or frozen, needs restarting or has a virus.
\end{chapteropening}

\section[la panne]{la panne — the breakdown, the failure}

\label{lesson:FR-C399-panne}



\begin{quote}
Before the new one: say the French for the link, then the French for to post.
\end{quote}

\subsection*{You'll want to know: la panne}
\textbf{la panne} — \textquotedblleft{}the breakdown, the failure\textquotedblright{}.

\begin{quote}
Ma voiture est en panne.
\end{quote}

\begin{quote}
\emph{My car has broken down.}
\end{quote}

\emph{Tomber en panne} — to break down. \emph{Une panne de courant} — a power cut.

\begin{cousinweb}
\emph{panne} was first a sailing word, \emph{mettre en panne}, for stopping a ship by turning its sails. The origin is not settled.
\end{cousinweb}

\begin{grammarlens}[title={What you've built}]
One more word for telling, explaining and describing at B1.
\end{grammarlens}

\subsection*{Your turn}
\begin{itemize}
\raggedright
  \item \emph{Say it:} \emph{la panne}
  \item \emph{Say it:} \emph{la panne}, once more
  \item \emph{Say it:} a sentence of your own with \emph{la panne}
  \item \emph{Recall:} say the French for the link, then the French for to post, then say \emph{la panne} again
\end{itemize}

\begin{tcolorbox}[breakable,colback=teal!4,colframe=teal!35!black,title={Before you move on}]
What does la panne mean? (\textquotedblleft{}The breakdown, the failure\textquotedblright{}.) Say it once more.
\end{tcolorbox}
\section[planter]{planter — to crash (of a computer, informal); to plant}

\label{lesson:FR-C399-planter}



\begin{quote}
Before the new one: say the French for to post, then the French for the breakdown.
\end{quote}

\subsection*{You'll want to know: planter}
\textbf{planter} — \textquotedblleft{}to crash (of a computer, informal); to plant\textquotedblright{}.

\begin{quote}
Mon ordinateur a planté pendant la réunion.
\end{quote}

\begin{quote}
\emph{My computer crashed during the meeting.}
\end{quote}

Informal in this sense. First, to plant: \emph{planter un arbre}.

\begin{cousinweb}
\emph{planter} is inherited from Latin \emph{plantāre} \textquotedblleft{}to plant\textquotedblright{} (EtymDB; Lewis \& Short), from \emph{planta} \textquotedblleft{}a sprout, a cutting\textquotedblright{}.
\end{cousinweb}

\begin{grammarlens}[title={What you've built}]
One more word for telling, explaining and describing at B1.
\end{grammarlens}

\subsection*{Your turn}
\begin{itemize}
\raggedright
  \item \emph{Say it:} \emph{planter}
  \item \emph{Say it:} \emph{planter}, once more
  \item \emph{Say it:} a sentence of your own with \emph{planter}
  \item \emph{Recall:} say the French for to post, then the French for the breakdown, then say \emph{planter} again
\end{itemize}

\begin{tcolorbox}[breakable,colback=teal!4,colframe=teal!35!black,title={Before you move on}]
What does planter mean? (\textquotedblleft{}To crash (of a computer, informal); to plant\textquotedblright{}.) Say it once more.
\end{tcolorbox}
\section[bloquer]{bloquer — to block; to freeze (reflexive, se bloquer)}

\label{lesson:FR-C399-bloquer}



\begin{quote}
Before the new one: say the French for the breakdown, then the French for to crash.
\end{quote}

\subsection*{You'll want to know: bloquer}
\textbf{bloquer} — \textquotedblleft{}to block; to freeze (reflexive, se bloquer)\textquotedblright{}.

\begin{quote}
L'écran s'est bloqué : rien ne bouge.
\end{quote}

\begin{quote}
\emph{The screen froze: nothing moves.}
\end{quote}

In Belgium, \emph{bloquer} also means to cram for exams.

\begin{cousinweb}
\emph{bloquer} is built on \emph{bloc} \textquotedblleft{}a block\textquotedblright{}, which came through Old French from Middle Dutch \emph{blok} \textquotedblleft{}a tree trunk, a log\textquotedblright{} (EtymDB).
\end{cousinweb}

\begin{grammarlens}[title={What you've built}]
One more word for telling, explaining and describing at B1.
\end{grammarlens}

\subsection*{Your turn}
\begin{itemize}
\raggedright
  \item \emph{Say it:} \emph{bloquer}
  \item \emph{Say it:} \emph{bloquer}, once more
  \item \emph{Say it:} a sentence of your own with \emph{bloquer}
  \item \emph{Recall:} say the French for the breakdown, then the French for to crash, then say \emph{bloquer} again
\end{itemize}

\begin{tcolorbox}[breakable,colback=teal!4,colframe=teal!35!black,title={Before you move on}]
What does bloquer mean? (\textquotedblleft{}To block; to freeze (reflexive, se bloquer)\textquotedblright{}.) Say it once more.
\end{tcolorbox}
\section[redémarrer]{redémarrer — to restart, to reboot}

\label{lesson:FR-C399-redemarrer}



\begin{quote}
Before the new one: say the French for to crash, then the French for to block.
\end{quote}

\subsection*{You'll want to know: redémarrer}
\textbf{redémarrer} — \textquotedblleft{}to restart, to reboot\textquotedblright{}.

\begin{quote}
Si ça ne marche pas, il faut redémarrer l'ordinateur.
\end{quote}

\begin{quote}
\emph{If it does not work, you have to restart the computer.}
\end{quote}

\emph{Re-} means again: to start again.

\begin{cousinweb}
\emph{redémarrer} is built from \emph{re-} and \emph{démarrer} \textquotedblleft{}to start\textquotedblright{}, first a sailing word for casting off a moored boat, from \emph{dé-} and \emph{amarrer} \textquotedblleft{}to moor\textquotedblright{}, a word borrowed from Dutch (EtymDB).
\end{cousinweb}

\begin{grammarlens}[title={What you've built}]
One more word for telling, explaining and describing at B1.
\end{grammarlens}

\subsection*{Your turn}
\begin{itemize}
\raggedright
  \item \emph{Say it:} \emph{redémarrer}
  \item \emph{Say it:} \emph{redémarrer}, once more
  \item \emph{Say it:} a sentence of your own with \emph{redémarrer}
  \item \emph{Recall:} say the French for to crash, then the French for to block, then say \emph{redémarrer} again
\end{itemize}

\begin{tcolorbox}[breakable,colback=teal!4,colframe=teal!35!black,title={Before you move on}]
What does redémarrer mean? (\textquotedblleft{}To restart, to reboot\textquotedblright{}.) Say it once more.
\end{tcolorbox}
\section[le virus]{le virus — the virus}

\label{lesson:FR-C399-virus}



\begin{quote}
Before the new one: say the French for to block, then the French for to restart.
\end{quote}

\subsection*{You'll want to know: le virus}
\textbf{le virus} — \textquotedblleft{}the virus\textquotedblright{}.

\begin{quote}
Si tu ouvres ce fichier, tu risques d'attraper un virus.
\end{quote}

\begin{quote}
\emph{If you open that file, you risk getting a virus.}
\end{quote}

For computers and for illness. The \emph{s} is said.

\begin{cousinweb}
\emph{virus} was taken whole from Latin \emph{vīrus} \textquotedblleft{}slime, poison\textquotedblright{} (EtymDB; Lewis \& Short); the computer sense came from English.
\end{cousinweb}

\begin{grammarlens}[title={What you've built}]
One more word for telling, explaining and describing at B1.
\end{grammarlens}

\subsection*{Your turn}
\begin{itemize}
\raggedright
  \item \emph{Say it:} \emph{le virus}
  \item \emph{Say it:} \emph{le virus}, once more
  \item \emph{Say it:} a sentence of your own with \emph{le virus}
  \item \emph{Recall:} say the French for to block, then the French for to restart, then say \emph{le virus} again
\end{itemize}

\begin{tcolorbox}[breakable,colback=teal!4,colframe=teal!35!black,title={Before you move on}]
What does le virus mean? (\textquotedblleft{}The virus\textquotedblright{}.) Say it once more.
\end{tcolorbox}
\section[(dialogue)]{Conditions et règles}

\label{lesson:FR-R399-revision-398-399}



\begin{quote}
Say the French for the network and the connection.
\end{quote}

\subsection*{Your turn}
Say these aloud:

\begin{itemize}
\raggedright
  \item the network, the connection, to browse, the link, to post
  \item the breakdown, to crash, to block, to restart, the virus
\end{itemize}

\begin{tcolorbox}[breakable,colback=teal!4,colframe=teal!35!black,title={Before you move on}]
The network? (\textbf{le réseau}.) The connection? (\textbf{la connexion}.) To browse? (\textbf{naviguer}.) The link? (\textbf{le lien}.) To post? (\textbf{publier}.) The breakdown? (\textbf{la panne}.) To crash? (\textbf{planter}.) To block? (\textbf{bloquer}.) To restart? (\textbf{redémarrer}.) The virus? (\textbf{le virus}.)
\end{tcolorbox}
\section[(dialogue)]{Ça ne marche pas — once more}

\label{lesson:FR-R399-revision-final-399}



\begin{quote}
Say the French for the breakdown and to crash.
\end{quote}

\subsection*{Your turn}
\begin{itemize}
\raggedright
  \item \emph{Say it:} the breakdown, to crash, to block, to restart, the virus
\end{itemize}

\begin{tcolorbox}[breakable,colback=teal!4,colframe=teal!35!black,title={Before you move on}]
The breakdown? (\textbf{la panne}.) To crash? (\textbf{planter}.) To block? (\textbf{bloquer}.) To restart? (\textbf{redémarrer}.) The virus? (\textbf{le virus}.)
\end{tcolorbox}
